\section{Markov property for simple SCMs}

By making use of the `acyclification', we can extend the global Markov property for L-CBNs to a global Markov property for simple SCMs.
The difference is that the Markov property for SCMs is formulated in terms of $\sigma$-separation rather than $d$-separation.
With the help of this Markov property, we can derive a very analogous theory for simple SCMs as for L-CBNs, with a do-calculus and adjustment.

\subsection{Acyclifications}

In Section~\ref{sec:graph_acyclifications}, we defined acyclifications of a CDMG.
We can also define an operation with the same name on SCMs.
\begin{Def}[Acyclification of SCM]\label{def:scm_acyclification}
  Let $M = \lt \Sin, \Send, \Srnd, \CX, P, f \rt$ be a simple SCM with graph $G(M)$.
  For each strongly connected component $C$ of $G(M)$ (i.e., a set of the form
  $\Sc^{G(M)}(v)$ for $v \in \Send$), let 
  $g_C : \CX_{\Sin \cup (\Send \setminus C) \cup \Srnd} \to \CX_C$
  be the unique solution function for $M^{[C]}$.
  Define $\tilde{f} : \CXJ \times \CXV \times \CXW \to \CXV$ by its components
  $$\tilde{f}_v(\xJ,\xV,\xW) = (g_{\Sc^{G(M)}(v)})_v(\xJ,x_{\Send\setminus C},\xW)$$
  for $v \in V$.
  $M^\acy = \ltuple \Sin, \Send, \Srnd, \CX, P, \tilde{f} \rtuple$ is called the
  \emph{acyclification of $M$}.
\end{Def}
The crucial property of this definition is the following result, which also motivates its name.
\begin{Prp}\label{prp:scm_acyclification_properties}
  Let $M = \lt \Sin, \Send, \Srnd, \CX, P, f \rt$ be a simple SCM. Its acyclification $M^\acy$ is
  acyclic and observably equivalent to $M$.
\end{Prp}
\begin{proof}
  We construct a directed graph $S$ from $G := G(M)$ with its strongly connected components $\{\Sc^G(v) : v \in \Send \cup \Sin \cup \Srnd\}$
  as nodes, and directed edges
  $C \tuh D$ if there is a directed edge $c \tuh d$ in $G$ with $c \in C, d \in D$ and $C \ne D$.
  The graph $S$ cannot contain a directed cycle, as that would imply the existence of a directed
  cycle in $G$ that traverses more than one of its strongly connected components. Hence $S$ is a DAG.
  
  Choose a topological ordering $<$ of $S$.
  Any node $C$ in $S$ can only have incoming directed edges
  in $S$ from $\Pred^S_<(C)$. This implies that for $v \in V$, $C = \Sc^G(v)$ can only have incoming
  edges in $G$ from $\bigcup \Pred^S_\le(C)$. That implies that the causal mechanism $f_C$
  can only depend on variables in $\bigcup \Pred^S_\le(C)$, and hence
  the unique solution function $g_C$, and therefore $\tilde{f}_C$, can depend on variables in $\bigcup \Pred^S_<(C)$ only.
  Therefore, for $v,w \in V$, a directed edge $w \tuh v$ in $G(M^\acy)$ implies 
  $w \in \bigcup \Pred^S_<(\Sc^G(v))$. We can therefore refine the topological ordering $<$ of $S$ 
  to a topological ordering of $G(M^\acy)$, by arbitrarily ordering the nodes within each strongly connected component
  of $G$. Hence $G(M^\acy)$ is acyclic.

  $M$ and $M^\acy$ are observably equivalent by construction: for all $x \in \CX$,
  \begin{align*}
  & x = \tilde{f}(x) \\
  & \iff \forall C \in S \cap V: x_C = \tilde{f}_C(x) \\
  & \iff \forall C \in S \cap V: x_C = g_C(x_J,x_{V\setminus C},x_W) \\
  & \iff \forall C \in S \cap V: x_C = f_C(x) \\
  & \iff x = f(x).
  \end{align*}
\end{proof}
The notion of acyclification of an SCM is compatible with that of a graph:
\begin{Prp}\label{prp:acyclifications_compatible}
  Let $M$ be a simple SCM.
  Then $G(M^\acy) \subseteq G'$ for any acyclification $G'$ of $G(M)$.
\end{Prp}
\begin{proof}
  By definition, the two graphs have the same nodes (input nodes $J$ and output nodes
  $V \cup W$). $G(M^\acy)$ has no bidirected edges, but $G'$ might. If there is a 
  directed edge $i \tuh j$ in $G(M^\acy)$ with $i \in J \cup V \cup W$ and $j \in V$, then
  the solution function $g_{\Sc^G(j)}$ of $M^{[\Sc^G(j)]}$ depends on $x_i$. 
  This can only happen if
  $i \notin \Sc^G(j)$ and $i$ is a parent of some $k$ according to $M$ with $k \in \Sc^G(j)$,
  i.e., if $i \tuh k$ in $G(M)$. In that case, $i \tuh j$ in $G'$ by definition of
  the graphical acyclification.
\end{proof}
Hence, two nodes in the same strongly connected component of $G(M)$ do not have any edge between them
in $G(M^\acy)$, whereas they necessarily have a connecting edge in any acyclification $G'$ of $G(M)$.
For two nodes in different strongly connected components of $G(M)$, the edges in $G(M^\acy)$ are also
present in $G'$, but not necessarily vice versa, as some parent-relations may cancel out in the solution function.

\subsection{Global Markov property for simple SCMs}

With the help of the acyclifications,
we can easily derive a Markov property for simple SCMs from the Markov property for CBNs by reducing
the general cyclic case to an acyclic case.
\begin{Cor}[Global Markov property for simple SCMs]\label{cor:gmp_simple_scm}
  Let $M = \lt \Sin, \Send, \Srnd, \CX, P, f \rt$ be a simple SCM with graph $G(M)$
  and Markov kernel $P_{M}(X_V,X_W \given \doit(X_J))$.
  Then for all $A, B, C \ins J \cup V \cup W$ (not necessarily disjoint) we have the implication:
  $$A \Perp^\sigma_{G(M)} B \given C \implies X_A \Indep_{P_{M}(X_V,X_W \given \doit(\XJ))} X_B \given X_C.$$
  If one wants to make the implicit dependence on $J$ more explicit one can equivalently also write:
  $$A \Perp^\sigma_{G(M)} J \cup B \given C \implies X_A \Indep_{P_{M}(X_V,X_W \given \doit(\XJ))} X_J,X_B \given X_C.$$
\end{Cor}
\begin{proof}
  Choose an acyclification $G'$ of $G(M)$. Then:
  \begin{equation*}\begin{split}
    A \Perp^\sigma_{G(M)} B \given C 
    & \iff A \Perp^d_{G'} B \given C \\
    & \implies A \Perp^d_{G(M^\acy)} B \given C \\
    & \implies X_A \Indep_{P_{M^\acy}(X_V,X_W \given \doit(\XJ))} X_B \given X_C \\
    & \iff X_A \Indep_{P_{M}(X_V,X_W \given \doit(\XJ))} X_B \given X_C.
  \end{split}\end{equation*}
  For the various implications / equivalences, we used:
  \begin{enumerate}
    \item $G'$ is an acyclification of $G(M)$ together with Proposition~\ref{prp:sigma_sep_as_d_sep}; 
    \item $G(M^\acy) \subseteq G'$ from Proposition~\ref{prp:acyclifications_compatible}, and that removing edges cannot turn a $d$-separation into a $d$-connection; 
    \item the global Markov property Theorem~\ref{thm-gmp-mCBN} for $M^\acy$ interpreted as a causal Bayesian network as in the proof of Proposition~\ref{prp:interventional_equivalence_cbn_scm} point \ref{prp:interventional_equivalence_cbn_scm:scm2cbn} (with deterministic Markov kernels for the endogenous variables, and purely probabilistic Markov kernels for the exogenous random variables), exploiting Proposition~\ref{prp:scm_acyclification_properties} that states that the acyclification $M^\acy$ is acyclic;
    \item by Proposition~\ref{prp:scm_acyclification_properties}, the acyclification $M^\acy$ has the same Markov kernel as the original SCM $M$.
  \end{enumerate}
\end{proof}
%This Markov property is not complete: it does not exploit all information in $G(M)$ that it possibly could to deduce 
%conditional independences. For example, endogenous variables without parents (for which the corresponding outcomes must be
%constant) could additionally be taken to block walks.
%\Joris{Should we change the condition for a non-collider to block the path if it is a deterministic function of $C$? We could detect this graphically by distinguishing deterministic kernels from non-deterministic ones (e.g.\ by a double circle).}

\subsection{Do-calculus for simple SCMs}\label{sec:scms_do-calculus}

%\Joris{Patrick also said that there are two subtleties one has to watch out for: (i) the `easy' condition to make
%the null sets work is to assume that each intervened Markov kernel appearing has a strictly positive
%density; (ii) the subtlety is that if one extends the output of the Markov kernel with a copy of (part of
%its) input, this assumption is already violated, so one should be quite careful when doing that.}

With the global Markov property for simple SCMs, it becomes straightforward to derive the do-calculus for simple SCMs.
First we will introduce some notation. The setting will be that a simple SCM $M = \lt \Sin, \Send, \Srnd, \CX, P, f \rt$
is given. 
As an auxiliary way to model hard interventions $\doit(B)$ for $B \subseteq V \cup W$ 
\JorisDoInput{if allowing intervention nodes for exogenous input nodes, $B \subseteq V \cup W \cup J$, the following needs to be updated; perhaps also don't do this inline but have a proper definition}
that leads to easy rules in the do-calculus, we introduce additional intervention variables
$I_b$ for $b \in B$. We denote $I_B := (I_b)_{b \in B}$.
This gives an extended SCM $\tilde{M} = \lt \Sin \cup I_B, \Send, \Srnd, \CX \times \prod_{b \in B} (\CX_b \dot\cup \{\star\}), P, \tilde{f} \rt$ with causal mechanism with components
$$\tilde{f}_b(x_{V \cup J \cup W},x_{I_B}) := \begin{cases}
  f_b(x_{V \cup J \cup W}) & x_{I_b} = \star \\
  x_{I_b} & x_{I_b} \in \CX_b
\end{cases}$$
for $b \in B \cap V$, and $\tilde{f}_v(x) := f_v(x_{V \cup J \cup W})$ for $v \in V \setminus B$.
Here, $x_{I_b} = \star$ encodes that there is no intervention on $b$, 
while $x_{I_b} \ne \star$ encodes that the hard intervention $\doit(X_b = x_{I_b})$ is performed.
We will denote the graph of $M$ by $G$ and the graph of $\tilde{M}$ by $G_{\doit(I_B)}$.
In the do-calculus, we make use of the extended graph $G_{\doit(I_B)}$ to test the separation statement, while the conclusion
about the properties of certain Markov kernels concerns those of the original SCM $M$.

%Remember the notation of Markov kernels for simple SCMs from Definition~\ref{def:simple_scm_Markov_kernels}.
The Markov kernel of simple SCM $M$ exists, is unique, and is denoted by $P_{M}(X_{V\cup W} \mid \doit(X_J))$.
For a subset $T \subseteq V \cup W$, we write 
$$P_{M}(X_{(V\cup W)\setminus T} \mid \doit(X_J,X_T)) := P_{M_{\doit(T)}}(X_{(V\cup W)\setminus T} \mid \doit(X_J,X_T)).$$
By conditioning on a subset $S \subseteq (V \cup W) \setminus T$, we obtain the conditional Markov kernel
$$P_{M}(X_{(V\cup W)\setminus (T \cup S)} \mid X_S, \doit(X_J,X_T)).$$

The only modification to the do-calculus for simple SCMs as compared to that for causal Bayesian networks is that
we have to replace all $d$-separations by $\sigma$-separations.
\begin{Thm}[Almost-sure do-calculus for simple SCMs, simplified]\label{thm:as-do-calculus-scms-simplified}
  Let $M = \lt \Sin, \Send, \Srnd, \CX, P, f \rt$ be a simple SCM with graph $G = G(M)$.
  Assume that we have $\sigma$-finite reference measures $\mu_v$ on $\Xcal_v$ for every $v \in V \cup W$
  and put $\mu_F := \bigotimes_{v \in F} \mu_v$ for $F \ins V \cup W$.
  Let $A,B,C \ins V \cup W$ and $D \ins V \cup W \cup J$ be such that $A,B,C,D$ are pairwise disjoint.
  Then we have the following 3 rules relating Markov kernels that can be generated from the SCM:
  \begin{enumerate}
      \item \emph{Insertion/deletion of observation}, for $J \subseteq D$: if
       \begin{align*} 
           A &\Perp^\sigma_{G_{\doit(D)}} B \given C \cup D, &
           \mu_{B \cup C} &\ll P_M(X_B,X_C|\doit(X_{D}))\ll \mu_{B \cup C},
       \end{align*}
        then:
          \[ P_M(X_A|X_B,X_C,\doit(X_{D})) = P_M(X_A|X_C,\doit(X_{D})) \qquad \mu_{B \cup C}\text{-a.s.}.\]
     \item \emph{Action/observation exchange}, for $J \subseteq D$: if
       \begin{align*} 
           A & \Perp^\sigma_{G_{\doit(I_B,D)}} I_B \given B \cup C \cup D, &          
                  \mu_{B \cup C} &\ll  P_M(X_B,X_C|\doit(X_{D})) \ll \mu_{B \cup C},\\
               &&   \mu_C &\ll  P_M(X_C|\doit(X_B,X_{D})) \ll \mu_C,
       \end{align*}
        then:
       \[ P_M(X_A|X_B,X_C,\doit(X_{D}))  = P_M(X_A|\doit(X_B),X_C,\doit(X_{D})) \qquad \mu_{B \cup C}\text{-a.s.}.\]
     \item \emph{Insertion/deletion of action}, for $J \subseteq D$: if
       \begin{align*} 
           A &\Perp^\sigma_{G_{\doit(I_B,D)}} I_B \given C  \cup D, &          
           \mu_C &\ll P_M(X_C|\doit(X_B,X_{D})) \ll \mu_C, \\
                 && \mu_C &\ll P_M(X_C|\doit(X_{D})) \ll \mu_C,
       \end{align*}
       then
       \[ P_M(X_A|\doit(X_B),X_C,\doit(X_{D})) = P_M(X_A|X_C,\doit(X_{D})) \qquad \mu_C\text{-a.s.}.\]
     \item \emph{Deletion of input}: 
         If 
         \begin{align*} 
           A &\Perp^\sigma_{G_{\doit(X_D)}} J \given C \cup D, &
             \mu_C &\ll P_M(X_C|\doit(X_{D\cup J})) \ll \mu_C,
         \end{align*}
         then there exists a Markov kernel $P_M(X_A|X_C,\doit(X_{D},\cancel{X_{J\sm D}}))$ such that:
         \[ P_M(X_A|X_C,\doit(X_{D},\cancel{X_{J\sm D}})) = P_M(X_A|X_C,\doit(X_{D \cup J})) \qquad \mu_C\text{-a.s.}. \]
    \end{enumerate}
\end{Thm}
\begin{proof}
  The proof is analogous to that of Corollary~\ref{cor:as-do-calculus-2}, except that it applies the global Markov property for simple SCMs, Corollary~\ref{cor:gmp_simple_scm}, instead of the one for causal Bayesian networks, Theorem~\ref{thm-gmp-mCBN}.
\end{proof}
\begin{comment}% OLD VERSION
\begin{Cor}[Do-calculus for simple SCMs]\label{cor:scm_do_calculus}
  Let $M = \lt \Sin, \Send, \Srnd, \CX, P, f \rt$ be a simple SCM with graph $G = G(M)$.
  Let $A,B,C \ins V \cup W$ and $D \ins V \cup W \cup J$ be such that $A,B,C,D$ are pairwise disjoint.
  \JorisDoInput{In rule 3, should we allow for $B \subseteq V \cup W \cup J$?}
  Then we have the following 3 rules relating marginal conditional intervened Markov kernels of $M$:
    \begin{enumerate}
        \item \emph{Insertion/deletion of observation}: If we have:
            \[ A \Perp^\sigma_{G_{\doit(D)}} B \given C \cup D,   \]
            then there exists a Markov kernel:
            \[ P\lp X_A|\cancel{X_B},X_C,\doit(X_{D\cup \cancel{J}})\rp:\;  \Xcal_C \times \Xcal_{D} \dshto \Xcal_A \]
            that is a version of:
            \[P_{M}\lp X_A|X_{B_2},X_C,\doit(X_{D\cup J})\rp,\] 
            for every subset $B_2 \ins B$ simultaneously.
            In short we could write:
            \[ P_{M}\lp X_A|X_B,X_C,\doit(X_{D\cup J})\rp = P_{M}\lp X_A|X_C,\doit(X_D)\rp.\]
            %\[ P\lp X_A|X_B,X_C,\doit(X_D)\rp =P\lp X_A|X_{B_2},X_C,\doit(X_D)\rp = P\lp X_A|X_C,\doit(X_D)\rp.\]
            Note that this Markov kernel is only dependent on $x_C$ and $x_{D}$, and constant in $x_{J\setminus D}$.
            Such a Markov kernel is unique up to $P_{M}\lp X_B,X_C|\doit(X_{D\cup J})\rp$-null sets.
            %then we get the equality:
            %\[ P\lp X_A|X_B,X_C,\doit(X_D)\rp = P\lp X_A|X_C,\doit(X_D)\rp,  \]
            %$P\lp X_B,X_C|\doit(X_D)\rp$-a.s., to be read as that there exists a version of the rhs 
            %that is also a version of the lhs. 
         \item \emph{Action/observation exchange}: If we have:
            \[ A \Perp^\sigma_{G_{\doit(I_B,D)}} I_B \given B \cup C \cup D,   \]
            then there exists a Markov kernel:
            \[P\lp X_A|\cancel{\doit}(X_B),X_C,\doit(X_{D\cup \cancel{J}})\rp: \;\Xcal_B \times \Xcal_C \times \Xcal_{D} \dshto \Xcal_A\]
            that is a version of:            
            \[ P_{M}\lp X_A|X_{B_1},\doit(X_{B_2}),X_C,\doit(X_{D\cup J})\rp, \]
            for every decomposition $B=B_1\dcup B_2$ simultaneously. In short we could write:
%            \[ P_{M}\lp X_A|\doit(X_B),X_C,\doit(X_{D\cup J})\rp  =  P_{M}\lp X_A|X_B,X_C,\doit(X_{D\cup J})\rp.\]
%            or even
            \[ P_{M}\lp X_A|\doit(X_B),X_C,\doit(X_D)\rp  =  P_{M}\lp X_A|X_B,X_C,\doit(X_D)\rp.\]
            since these Markov kernels are constant in $x_{J\setminus D}$.
            Such a Markov kernel is unique up to $P_{M}\lp X_B,X_C|\doit(X_{I_B},X_{D\cup J})\rp$-null sets.
         \item \emph{Insertion/deletion of action}: If we have:
            \[ A \Perp^\sigma_{G_{\doit(I_B,D)}} I_B \given C  \cup D,   \]
            then there exists a Markov kernel: 
            \[P\lp X_A|\cancel{\doit(X_B)},X_C,\doit(X_{D\cup \cancel{J}})\rp:\; \Xcal_C \times \Xcal_{D} \dshto \Xcal_A\] 
            that is a version of:  
            \[ P_{M}\lp X_A|\doit(X_{B_2}),X_C,\doit(X_{D\cup J})\rp\]
            for every subset $B_2 \ins B$ simultaneously. In short we could write:
            \[ P_{M}\lp X_A|\doit(X_B),X_C,\doit(X_{D\cup J})\rp = P_{M}\lp X_A|X_C,\doit(X_{D})\rp.\]
            Note that this Markov kernel is only dependent on $x_C$ and $x_{D}$, and constant in $x_{J\setminus D}$.
            Such a Markov kernel is unique up to $P_{M}\lp X_C|\doit(X_{I_B},X_{D\cup J})\rp$-null set.
%          \item \Joris{\emph{Deletion of action II}: Here we take $B \subseteq J$ disjoint from $A,C,D$.
%            If we have:
%            \[ A \Perp^\sigma_{G_{\doit(D)}} \emptyset \given C \cup D, \]
%            then there exists a Markov kernel: 
%            \[P\lp X_A|X_C,\doit(X_{D\cup \cancel{J}})\rp:\; \Xcal_C \times \Xcal_{D} \dshto \Xcal_A\] 
%            that is a version of:  
%            \[ P_{M}\lp X_A|X_C,\doit(X_{D\cup J})\rp\]
%            In short we could write:
%            \[ P_{M}\lp X_A|X_C,\doit(X_{D\cup J})\rp = P_{M}\lp X_A|X_C,\doit(X_{D})\rp.\]
%            Note that this Markov kernel is only dependent on $x_C$ and $x_{D}$, but constant in $x_{J\setminus D}$.
%            Such a Markov kernel is unique up to $P_{M}\lp X_C|\doit(X_{D\cup J})\rp$-null set.}
    \end{enumerate}
\end{Cor}
\begin{proof}
  The proof is also analogous to that of Corollary~\ref{cor:do-calculus}, except that it applies the global
  Markov property for simple SCMs, Corollary~\ref{cor:gmp_simple_scm}, instead of the one for causal Bayesian networks,
  Theorem~\ref{thm-gmp-mCBN}. 
  %Furthermore, we do not assume that $J \subseteq D$, so therefore in each Markov kernel
  %from the SCM (but not the ones whose existence is part of the conditional independence statement) we replaced 
  %$\doit(X_D)$ by $\doit(X_{D \cup J})$.

%\Joris{Rule 4. The assumption:
%            \[ A \Perp^\sigma_{G_{\doit(D)}} B \given C  \cup D,   \]
%implies the conditional independence by GMP \ref{cor:gmp_simple_scm}:
%  \[ X_A \Indep_{P_{M}(X_V,X_W|\doit(X_{D\cup J}))} X_B \given X_C,X_D.\]
%So we have the following factorization:
%  \[ P_{M}\lp X_A,X_C|\doit(X_{D\cup J})\rp = Q(X_A|X_C,X_D) \otimes P_{M}\lp X_C|\doit(X_{D\cup J})\rp,\]
%for some Markov kernel $Q(X_A|X_C,X_D)$, which serves as a version of the conditional Markov kernel:
%  \[P_{M}\lp X_A|X_C,\doit(X_{D\cup J})\rp,\]
%and which is constant in $X_{B}$ and in $X_{J\setminus D}$.
%So we could write
%  \[P_{M}\lp X_A|X_C,\doit(X_{D\cup J})\rp = P_{M}\lp X_A|X_C,\doit(X_D)\rp\]
%}
\end{proof}
\end{comment}

%Thus, by making use of the do-calculus, 
%we can in some cases derive the existence of Markov kernels $P_{M}(X_A \given \doit(X_B))$
%where $J \not\subseteq B$;
%these should be interpreted as Markov kernels $\CX_{J \cup B} \dshto \CX_{A}$ that are constant in 
%$x_{J\setminus B}$. 
%Essentially, the do-calculus follows immediately from application of the global Markov property. 
While the derivation of the do-calculus relies essentially on the global Markov property, sometimes one
can make use of the global Markov property and a more careful analysis of null sets to obtain stronger conclusions.
In particular, Proposition~\ref{prp:do-calculus} and Theorem~\ref{thm:as-do-calculus-1} also hold for simple SCMs
if one replaces the d-separation statements by the analogous $\sigma$-separation statements.

%\begin{Rem}
%  In case you are wondering how to use rule 3 to insert or delete an action for an exogenous input
%  variable $j \in J$: this can be done via the special case $B = I_B = \emptyset$ and $D = J \setminus \{j\}$.
%  \Joris{Not possible anymore. Workaround? Add rule 3b?}
%  \JorisDoInput{This would become redundant}
%\end{Rem}

%\Joris{
%\begin{Rem}
%I conjecture, for $B \subseteq V \cup W \cup J$ disjoint from $A \cup C \cup D$:
%\[ A \Perp^\sigma_{G_{\doit(I_{B\setminus J},D)}} J \given C \cup D \iff
%\forall B_2 \subseteq B \setminus J: A \Perp^\sigma_{G_{\doit(B_2,D)}} B \given C \cup D, \]
%\end{Rem}
%}

\subsection{Adjustment}\label{sec:scms_adjustment}

Let $M = \lt \Sin, \Send, \Srnd, \CX, P, f \rt$ be a simple SCM.
Let us assume $J \subseteq D$.
We are interested in estimating the conditional causal effect:
\[ P_{M}(X_A|X_C,\doit(X_B,X_D)), \]
but we only have data from:
\[ P_{M}(X_A,X_B,X_F|X_C,\doit(X_D)). \]
The following (pairwise disjoint) index sets will have the following roles:
\begin{description}
  \item[$A \ins V$]: the outcome variables of interest.
  \item[$B \ins V \cup W$]: the treatment or intervention variables.
  \item[$C \ins V \cup W$]: general conditional (context) variables under which the data was collected.
  \item[$J \ins D \ins V \cup W \cup J$]: general interventional (context) variables that were set by the experimenter.
  \item[$F_0 \ins V \cup W$]: core adjustment variables, i.e.\ features that were measured.
  \item[$F_1 \ins V \cup W$]: additional measured adjustment variables, with $F = F_0 \cup F_1$.
  \item[$H \ins V \cup W$]: additional unobserved variables.
\end{description}
We will make use of the same extended SCM $\tilde{M}$ with intervention variables $I_b$ for $b \in B$ and graph
$G_{\doit(I_B)}$ as for stating the do-calculus.

\begin{comment}
\begin{Ass}\label{ass:positivity}
  Let $M = \lt \Sin, \Send, \Srnd, \CX, P, f \rt$ be a simple SCM.
  Assume that:
  \begin{enumerate}
    \item Each measurable space $\Xcal_u$ for $u \in \Send \cup \Srnd$ comes equipped with a fixed measure $\mu_u$, and
    \item For every subset $D \ins \Send$ the intervened Markov kernel $P_M(X_D \mid \doit(X_{J \cup V \cup W \sm D}))$ is
      absolute continuous w.r.t.\ the product measure $\mu_D:=\bigotimes_{v \in D} \mu_v$ and vice versa:
            \[ \mu_D \ll P_M(X_D \mid \doit(X_{J \cup V \cup W \sm D})) \ll \mu_D. \]
  \end{enumerate}
%  The second point is satisfied if each intervened Markov kernel $P_M(X_u \mid \doit(X_{J \cup V \sm \{u\}}))$ has a strictly positive Doob-Radon-Nikodym derivative/density w.r.t.\ the measure $\mu_u$ for all $u \in \Send \cup \Srnd$.
%  \Joris{Is it? May not be due to cycles}
\end{Ass}
\end{comment}
\begin{Thm}[General adjustment formula for simple SCMs]\label{thm:gen-adjustment-scms}
Given a simple SCM $M = \lt \Sin, \Send, \Srnd, \CX, P, f \rt$ with graph $G = G(M)$. %that satisfies Assumption~\ref{ass:positivity},
Assume that all the following $\sigma$-separations hold in the graph $G_{\doit(I_B,D)}$:
    \begin{align}
        (F_0 \cup H) & \Perp^\sigma_{G_{\doit(I_B,D)}} I_B \given (C \cup D), \\
        A & \Perp^\sigma_{G_{\doit(I_B,D)}} (F_1 \cup I_B) \given (B \cup F_0 \cup H \cup C \cup D), \\
                      H &\Perp^\sigma_{G_{\doit(I_B,D)}} B \given (F \cup C \cup  I_B \cup D).
    \end{align}
Further assume that we have reference measures $\mu_v$ on $\Xcal_v$, $v \in V \cup H$, such that:
  \begin{align*}
      \mu_{B \cup C \cup F \cup H} &\ll  P(X_B,X_C,X_F,X_H|\doit(X_D)) \ll \mu_{B \cup C \cup F \cup H},\\
      \mu_{C \cup F \cup H} &\ll  P(X_C,X_F,X_H|\doit(X_B,X_D)) \ll \mu_{C \cup F \cup H}.
  \end{align*}
Then we have the adjustment formula:
\[ P_{M}(X_A|X_C,\doit(X_B,X_D)) = P_{M}(X_A|X_B,X_C,X_F,\doit(X_D)) \circ P_{M}(X_F|X_C,\doit(X_D)) \qquad \mu_{B \cup C}\text{-a.s.} \]
\end{Thm}
\begin{proof}
  Analogous to that of Theorem~\ref{thm-gen-adjustment}, but now using
  the global Markov property for simple SCMs, Corollary~\ref{cor:gmp_simple_scm}, 
  instead of the one for causal Bayesian networks, Theorem~\ref{thm-gmp-mCBN}. 
\end{proof}
%Note that the ``a.s.'' qualifier is imprecise, a more precise statement could be obtained by tracing the uniqueness of the various Markov kernels appearing in the proof.
The following is just the special case $F_1=H=\emptyset$.
\begin{Cor}[Conditional interventional backdoor covariate adjustment formula for simple SCMs]\label{cor:backdoor-cond-int-scms}
  Let $M = \lt \Sin, \Send, \Srnd, \CX, P, f \rt$ be a simple SCM with graph $G = G(M)$. % that satisfies Assumption~\ref{ass:positivity}.
  Assume that the \emph{conditional interventional backdoor criterion} in the graph $G_{\doit(I_B,D)}$ holds:
  \begin{enumerate}
    \item $\displaystyle F \Perp^\sigma_{G_{\doit(I_B,D)}} I_B  \given (C \cup D)$, and:
    \item $\displaystyle A \Perp^\sigma_{G_{\doit(I_B,D)}} I_B  \given (B \cup F \cup C \cup D)$.
  \end{enumerate}
  Further assume the following absolute continuities:
  \begin{align*}
      \mu_{B \cup C \cup F} &\ll  P(X_B,X_C,X_F|\doit(X_D)) \ll \mu_{B \cup C \cup F},\\
      \mu_{C \cup F} &\ll  P(X_C,X_F|\doit(X_B,X_D)) \ll \mu_{C \cup F}.
  \end{align*}
  Then we have the adjustment formula:
  \[ P_{M}(X_A|X_C,\doit(X_B,X_D)) = P_{M}(X_A|X_B,X_C,X_F,\doit(X_D)) \circ P_{M}(X_F|X_C,\doit(X_D)) \qquad \mu_{B \cup C}\text{-a.s.} \]
\end{Cor}

Without the conditioning set, i.e.\ $C=\emptyset$, and direct careful analysis we get a version with slightly weaker positivity assumptions:
\begin{Thm}[Interventional backdoor covariate adjustment formula]\label{thm:backdoor-int-scms}
  Let $M = \lt \Sin, \Send, \Srnd, \CX, P, f \rt$ be a simple SCM with graph $G = G(M)$. % that satisfies Assumption~\ref{ass:positivity}.
  Assume that the \emph{interventional backdoor criterion} in the graph $G_{\doit(I_B,D)}$ holds:
  \begin{enumerate}
    \item $\displaystyle F \Perp^\sigma_{G_{\doit(I_B,D)}} I_B  \given D$, and:
    \item $\displaystyle A \Perp^\sigma_{G_{\doit(I_B,D)}} I_B  \given (B \cup F \cup D)$.
  \end{enumerate}
  Further assume the following absolute continuity:
  \begin{align*}
      P_M(X_F|\doit(X_D)) \otimes P_M(X_B|\doit(X_D)) &\ll  P_M(X_F,X_B|\doit(X_D)).
  \end{align*}
  Then we have the adjustment formulas: 
  \begin{align*}
    P_M(X_A,X_F|\doit(X_B,X_D)) &= P_M(X_A|X_F,X_B,\doit(X_D)) \otimes P_M(X_F|\doit(X_D)) && P_M(X_B|\doit(X_D))\text{-a.s.}, \\
    P_M(X_A|\doit(X_B,X_D)) &= P_M(X_A|X_F,X_B,\doit(X_D)) \circ P_M(X_F|\doit(X_D)) && P_M(X_B|\doit(X_D))\text{-a.s.}
  \end{align*}
\end{Thm}
\begin{proof}
  Analogous to that of Corollary~\ref{cor:backdoor-int}, but now using
  the global Markov property for simple SCMs, Corollary~\ref{cor:gmp_simple_scm}, 
  instead of the one for causal Bayesian networks, Theorem~\ref{thm-gmp-mCBN}. 
\end{proof}
We can now further specialize to the case with $C=D=J= \emptyset$ and immediately get:
\begin{Cor}[Backdoor covariate adjustment for simple SCMs]\label{cor:backdoor-scms}
  Given a simple SCM $M = \lt \Sin, \Send, \Srnd, \CX, P, f \rt$ with graph $G = G(M)$. % that satisfies Assumption~\ref{ass:positivity},
  Assume that the \emph{backdoor criterion} holds:
  \begin{enumerate}
    \item $\displaystyle F \Perp^\sigma_{G_{\doit(I_B)}} I_B$, and:
    \item $\displaystyle A \Perp^\sigma_{G_{\doit(I_B)}} I_B  \given (B \cup F)$.
  \end{enumerate}
  Further assume the following absolute continuity:
  \begin{align*}
      P_M(X_F) \otimes P_M(X_B) &\ll  P_M(X_F,X_B).
  \end{align*}
%  Then we have the adjustment formula:
%  \[ P_{M}(X_A|\doit(X_B)) = P_{M}(X_A|X_B,X_F) \circ P_{M}(X_F) \qquad \text{$\mu_{V\cup W}$-a.s.} \]
  Then we have the adjustment formulae: 
  \begin{align*} 
    P_M(X_A,X_F|\doit(X_B)) &= P_M(X_A|X_F,X_B) \otimes P_M(X_F)  && P_M(X_B)\text{-a.s.}, \\
    P_M(X_A|\doit(X_B)) &= P_M(X_A|X_F,X_B) \circ P_M(X_F)   && P_M(X_B)\text{-a.s.}
  \end{align*}
\end{Cor}
%\Joris{The positivity conditions are pretty strong. Cannot we throw them overboard by 
%showing that
%\[ P_{M}(X_A|\doit(X_B)) = 
%P_{M}(X_A|X_B,X_F) \circ P_{M}(X_F) \qquad \text{$P(X_F,X_B)$-a.s. ?}
%\]
%Patrick notes that the latter thing is not really formally defined for the Markov kernel on the l.h.s..
%I countered that perhaps we should then define this notion, perhaps just by extending the one on the l.h.s.\
%by making it constant in $X_F$, i.e., as $P_{M}(X_A|\doit(X_B,\cancel{X_F}))$.
%At closer inspection, it seems that this is not such a good idea at all.
%However, what we could do in practice (in the discrete case) is to propagate bounds for anything that cannot be estimated from the data.
%}
The literature often fails to mention the strict positivity assumptions, even though without sufficient positivity, the various backdoor criteria may not hold.
A simple example of how the adjustment formula may fail if the strict positivity assumptions are not met is provided in Example~\ref{eg:counterexample_nonpositive_backdoor_adjustment}.


\subsection{Bounds on causal effects}

If causal effects cannot be identified from the observable distribution, we may still be able to derive informative bounds (see also Figure~\ref{fig:natural_bounds}).
We first prove the \emph{consistency} property of solution functions of simple SCMs.
\begin{Prp}\label{prp:simple_scm_consistency}
  Let $M = \lt \Sin, \Send, \Srnd, \CX, P, f \rt$ be a simple SCM.
  Let $g : \CXJ \times \CXW \to \CXV$ be the solution function of $M$.
  Let $V = A \dcup B$ be a partition of the endogenous variables of $M$,
  and let $\tilde{g} : \CX_A \times \CXJ \times \CXW \to \CX_B$ be the solution function of $M_{\doit(X_A)}$. Then
  $$g_B(x_J,x_W) = \tilde{g}(g_A(x_J,x_W),x_J,x_W)$$
  for all $x \in \CX$.
\end{Prp}
\begin{proof}
  For all $x \in \CX$:
  $$\begin{cases}
    x_A & = f_A(x), \\
    x_B & = f_B(x)
  \end{cases} \iff
  \begin{cases}
    x_A & = g_A(x_J,x_W), \\
    x_B & = g_B(x_J,x_W).
  \end{cases}$$
  Also, for all $x \in \CX$:
  $$x_B = f_B(x)
   \iff x_B = \tilde{g}_B(x_A,x_J,x_W).$$
  Hence, for all $x \in \CX$:
  $$\begin{cases}
    x_A & = f_A(x), \\
    x_B & = f_B(x)
  \end{cases} \implies
  \begin{cases}
    x_A & = g_A(x_J,x_W), \\
    x_B & = \tilde{g}_B(x_A,x_J,x_W)
  \end{cases}\implies
  \begin{cases}
    x_A & = g_A(x_J,x_W), \\
    x_B & = \tilde{g}_B(g_A(x_J,x_W),x_J,x_W).
  \end{cases}$$
  The uniqueness of the solution function $g$ of $M$ now implies the consistency statement.
\end{proof}

\begin{figure}
  \begin{center}
    \begin{tikzpicture}
      \begin{scope}[xshift=0]
        \node[ndout] (X1) at (0,0) {$X_1$};
        \node[ndout] (X2) at (3,0) {$X_2$};
        \node[ndout] (X3) at (1.5,2.5) {$X_3$};
%        \node[ndlat] (E1) at (0,1.5) {$X_A$};
        \draw[arout] (X1) -- (X2);
        \draw[arout] (X3) -- (X1);
        \draw[arout] (X3) -- (X2);
        \node at (-1,0) {(a)};
      \end{scope}
      \begin{scope}[xshift=6cm]
        \node[ndout] (X1) at (0,0) {$X_1$};
        \node[ndout] (X2) at (3,0) {$X_2$};
        \node[ndout] (X3) at (1.5,2.5) {$X_3$};
%        \node[ndlat] (E1) at (0,1.5) {$X_A$};
        \draw[arout] (X1) -- (X2);
        \draw[arout, bend left] (X2) edge (X1);
        \draw[arout] (X3) -- (X1);
        \draw[arout] (X3) -- (X2);
        \draw[arout, bend left] (X1) edge (X3);
        \draw[arout, bend right] (X2) edge (X3);
        \draw[huh, bend right] (X2) edge (X1);
        \draw[huh, bend left] (X3) edge (X1);
        \draw[huh, bend left] (X2) edge (X3);
        \node at (-1,0) {(b)};
      \end{scope}
    \end{tikzpicture}
  \end{center}
  \caption{Two graphs of simple SCMs. In (a), the Markov kernel $P_M(X_2 \given \doit(X_1), X_3)$ is identifiable (under positivity assumptions) from $P_M(X_1,X_2,X_3)$. In (b), it is not identifiable, but we can still bound it using the natural bounds if both $X_1$ and $X_3$ are discrete.\label{fig:natural_bounds}}
\end{figure}

\cite{ManskiNagin1998} proved the following `natural' bounds. We point out here that they also hold for simple SCMs.
\begin{Thm}[Natural bounds on causal effect]\label{thm:natural_bounds}
  Let $M$ be a simple SCM with endogenous variables $V \supseteq \{1,2,3\}$ and no exogenous input variables.
  Assume that $\CX_1$ is discrete.
  Then
  \begin{equation}\label{eq:natural_bounds}\begin{split}
    P_M(X_1=a,X_2 \in B|X_3 \in C) &\le P_M(X_2\in B|X_3 \in C, \doit(X_1=a)) \\
                   &\le P_M(X_1=a,X_2\in B|X_3 \in C) + P_M(X_1\ne a|X_3 \in C).
  \end{split}\end{equation}
  for any $a\in\C{X}_1$, measurable $B \subseteq \C{X}_2$, and measurable $C \ins \C{X}_3$ with $P_M(X_3 \in C) > 0$.
\end{Thm}
\begin{proof}
  We use the consistency $X_2=X_2^{\doit(X_1=X_1)}$ and elementary probability theory:
    \begin{align*}
      P_M(X_1=a,X_2\in B|X_3 \in C) 
      &= P_M(X_1=a,X_2^{\doit(X_1=a)}\in B|X_3 \in C) \\
      &\le P_M(X_2^{\doit(X_1=a)}\in B|X_3 \in C) \\
      &= P_M(X_1=a,X_2^{\doit(X_1=a)}\in B|X_3 \in C) \\
      & \quad + P_M(X_1\ne a,X_2^{\doit(X_2=a)}\in B|X_3 \in C) \\
      &\le P_M(X_1=a,X_2\in B|X_3 \in C) + P_M(X_1 \ne a|X_3 \in C)
    \end{align*}
  The statement follows since
  $$ P_M(X_2\in B|X_3 \in C, \doit(X_1=a)) = P_M(X_2^{\doit(X_1=a)}\in B|X_3 \in C).$$
\end{proof}
This so-called ``natural'' bound can be shown to be tight. 
Remarkably, we do not need to make any assumptions regarding the causal relations between the three endogenous variables.
This allows us to bound the causal effect of $X_1$ on $X_2$ in the presence of confounding and cycles.
Unfortunately, it can be shown that there exists no analogous bound in case $X_1$ is real-valued.

If one has a priori knowledge about the range of $X_2$ (in case it is real-valued), one can also derive a bound on the expected result of an intervention \cite{ManskiPepper2013}.
\begin{Cor}\label{cor:natural_bounds}
  In the situation of Theorem~\ref{thm:natural_bounds}, suppose that $\CX_2 = [\alpha,\beta]$ and $0 < P_M(X_1=a|X_3\in C) < 1$. Then
  \begin{equation}\label{eq:natural_bounds_CATE}\begin{split}
    & P_M(X_1=a|X_3\in C) \Exp_M(X_2|X_1=a,X_3\in C) + \alpha P_M(X_1\ne a|X_3\in C) \\
    {}\le {}& \Exp_M(X_2|X_3\in C,\doit(X_1=a)) \\
    {}\le {}& P_M(X_1=a|X_3\in C) \Exp_M(X_2|X_1=a,X_3\in C) + \beta P_M(X_1\ne a|X_3\in C).
  \end{split}\end{equation}
\end{Cor}
\begin{proof}
Using consistency, we get:
\begin{equation*}\begin{split}
  &P_M(X_2^{\doit(X_1=a)} \in B|X_3\in C) \\
  &= P_M(X_1=a|X_3\in C)  P_M(X_2^{\doit(X_1=a)} \in B|X_1=a,X_3\in C)  \\
  & + P_M(X_1 \ne a|X_3\in C) P_M(X_2^{\doit(X_1=a)} \in B|X_1\ne a,X_3\in C) \\
  &= P_M(X_1=a|X_3\in C)  P_M(X_2 \in B|X_1=a,X_3\in C)  \\
  & + P_M(X_1 \ne a|X_3\in C) P_M(X_2^{\doit(X_1=a)} \in B|X_1\ne a,X_3\in C).
\end{split}\end{equation*}
Integrating over $X_2$, the assumption $\CX_2 = [\alpha,\beta]$, interval arithmetic, and affinity of expected values, gives:
\begin{equation*}\begin{split}
  &\Exp_M(X_2^{\doit(X_1=a)}|X_3\in C) \\
  &= P_M(X_1=a|X_3\in C) \Exp_M(X_2|X_1=a,X_3\in C)  \\
  &+ P_M(X_1 \ne a|X_3\in C) \Exp_M(X_2^{\doit(X_1=a)}|X_1\ne a,X_3\in C) \\
  &\in P_M(X_1=a|X_3\in C) \Exp_M(X_2|X_1=a,X_3\in C) + P_M(X_1 \ne a|X_3\in C) [\alpha,\beta].
\end{split}\end{equation*}
\end{proof}
